\documentclass[10pt,a4paper]{amsart}
\usepackage[margin=19mm]{geometry}
\usepackage{amsmath,amssymb,mathtools}
\usepackage[scr=rsfs,cal=euler]{mathalfa}
\usepackage{xfrac}
\usepackage[unicode,hidelinks,bookmarks=false]{hyperref}
\newcommand{\ie}{i.e.\ }
\newcommand{\cf}{cf.\ }
\newcommand{\resp}{respectively\ }
\newcommand{\Subsection}{\subsection}
\providecommand{\nobreakdash}{\nobreak\hbox{-}}
\setlength{\emergencystretch}{2em}
\multlinegap0pt
\usepackage{amssymb}
\usepackage[shortlabels]{enumitem}
\usepackage{dsfont}
\usepackage{float}
\usepackage{xcolor}
\newtheorem{theorem}{Theorem}
\newcommand{\dik}{\text{{\rm{DIK}}}}
\newcommand{\jac}{\text{{\rm{Jac}}}}
\title{Weighted averages of $p$-adic hypergeometric functions and traces of Frobenius of elliptic curves}
\author{Riya Mandal, Neelam Saikia}
\date{Source: arXiv:2603.01148v2 (CC BY 4.0)}
\begin{document}
\maketitle

\noindent\emph{Source and licence.} Weighted averages of $p$-adic hypergeometric functions and traces of Frobenius of elliptic curves. Version arXiv:2603.01148v2 (CC BY 4.0), distributed by arXiv under the Creative Commons Attribution 4.0 International licence, http://creativecommons.org/licenses/by/4.0/, as recorded in the arXiv licence metadata for this exact version and shown on the abstract page https://arxiv.org/abs/2603.01148v2. Full licence notice: \texttt{licenses/arxiv-2603.01148-CC-BY-4.0.txt}.\par\smallskip

\noindent\textit{Editorial scope.} Counted unit: the article's theorem theorem-Jacobi (source0.tex lines 215--223). The article's complete original proof of it is delivered with it (source0.tex lines 578--623, one \texttt{proof} environment of 46 lines, which the article places away from the statement). No mathematics is written, repaired or re-derived: the statement and the proof are the article's own lines, in the article's own notation, and no retained line is substituted or re-flowed.  Retained uncounted as the source-local context the proof needs: lines 176--182; lines 298--301; lines 364--366.  Nothing was omitted from any retained span, so the package carries no omission record.  No retained line contains a citation, so the wrapper carries no bibliography and no bibliography entry.  No retained line contains a figure environment, an image-inclusion command or drawing code, so no image is delivered and the package declares no asset dependency.  Editorial formatting only, in addition: an AMS \texttt{amsart} preamble that restates the article's own declarations and macros used by the retained lines, namely the environments \texttt{theorem} on separate counters, together with the standard packages those lines need.  The retained environments print in this document as Theorem 1 in order of appearance, whereas the article prints the same items with the numbering of its own full document.  The complete native source of the article as distributed in the arXiv e-print for this exact version is bundled unchanged as \texttt{sources/195545/source0.tex}.
	\begin{theorem}\label{theorem-DIK-1}
For a prime $p>3$ and $\lambda\neq0,\frac{9}{4}$ we have 
	\begin{align}
\frac{p^2\phi(-3\lambda)}{(p-1)} \sum_{r=1}^{p-2}{\omega}^r(6\lambda)  \binom{\overline{\omega}^r}{\overline{\omega}^r\phi}
{_2G_2}(r,1)_p=\frac{\phi(3\lambda)}{(p-1)} + a_p(E_{\lambda}^{\dik}).\notag
   \end{align}
	\end{theorem}

\begin{equation}\label{inverse}
    g(\chi)g(\overline{\chi})=
        p\cdot\chi(-1). 
        \end{equation}

\begin{align}\label{prod-3}
   \Gamma_p\left(\left\langle\frac{tj}{p-1}\right\rangle\right)\omega(t^{tj}) \prod\limits_{h=1}^{t-1}\Gamma_p\left(\frac{h}{t}\right)=\prod\limits_{h=0}^{t-1}\Gamma_p\left(\left\langle\frac{h}{t}+\frac{j}{p-1}\right\rangle\right).
\end{align}

	\begin{theorem}\label{theorem-Jacobi}
	Let $p>3$ be a prime and $\lambda\in \mathbb{F}_p$ such that $\lambda\neq0, \pm 1.$ Then we have 
	$$ 
	a_p(E^{\jac}_{\lambda})=1+\phi(2\lambda) \cdot{_2G_2}\left[\begin{array}{cc}
\frac{1}{4}   & \frac{3}{4} \\
  0    & 0  \end{array}|{\lambda^2}\right]_p=1+p\cdot\phi(-2\lambda)\cdot {_2G_2}\left[\begin{array}{cc}
\frac{1}{4}   & \frac{3}{4} \\
  \frac{1}{2}    & \frac{1}{2}  \end{array}|{\lambda^2}\right]_p.$$
\end{theorem}

\begin{proof}[Proof of Theorem \ref{theorem-Jacobi}] We first prove the first equality. To do this, we 
employ similar method as use in proof of Theorem \ref{theorem-DIK-1} and derive that
\begin{align}
 -p(a_p(E_\lambda^{\jac}))&=p+\frac{g(\phi)\phi(-1)}{(p-1)^3}\sum_{r,s,t=0}^{p-2}g(T^{-r})g(T^{-s})g(T^{-t})T^s(2\lambda)\sum_{z}\phi T^{r+s+t}(z)\sum_{x}T^{4r+2s}(x).\notag
\end{align}
Using orthogonality of multiplicative characters it is easy to see that the last sum present in the above identity is nonzero if  $s=-2r$ and $s=-2r+\frac{p-1}{2}.$ Substituting these values in the above identity we have 
\begin{align}\label{eqn-1004}
-p\cdot a_p(E^{\jac}_{\lambda})&=p+A_1+A_2,
\end{align}
\noindent where $$A_1=\frac{g(\phi)\phi(-1)}{(p-1)^2}\sum\limits_{r,t=0}^{p-2}g(T^{-r})g(T^{2r})g(T^{-t})T^{-2r}(2\lambda)\sum_{z}T^{t-r}(z)\phi(z)$$ and 
where $$A_2=\frac{g(\phi)\phi(-2\lambda)}{(p-1)^2}\sum_{r,t=0}^{p-2}g(T^{-r})g(T^{2r}\phi)g(T^{-t})T^{-2r}(2\lambda)\sum_{z\in \mathbb F_p^*}T^{t-r}(z).$$
Now, using orthogonality relation of multiplicative characters and then applying Hasse-Davenport relation in the expression of $A_1,$ we deduce that $$A_1=-2p.$$
\noindent Similarly, using orthogonality relation of multiplicative characters we obtain
\begin{align}\label{eqn-1005}
  A_2&=\frac{1}{(p-1)}\sum_{r=0}^{p-2}g(T^{-r})g(T^{2r}\phi)g(T^{-r})g(\phi)\phi(-2\lambda)T^{r}\left(\frac{1}{4\lambda^2}\right).
\end{align}
Again, by making use of Hasse-Davenport formula we have $$g(T^{2r}\phi)=\frac{g(T^{4r}){T}^{-4r}(2)g(\phi)}{g(T^{2r})}.$$
Substituting this into \eqref{eqn-1005} and using \eqref{inverse} we have 
\begin{align}\label{eqn-1006}
  A_2&=\frac{p\phi(2\lambda)}{(p-1)}\sum_{r=0}^{p-2}g(T^{-r})^2\frac{g(T^{4r})}{g(T^{2r})}T^{r}\left(\frac{1}{2^6\lambda^2}\right).
\end{align}
Now, replacing $T$ by $\overline{\omega}$ and applying Gross-Koblitz formula we have 
\begin{align}\label{eqn-1010}
A_2=\frac{p \phi(2\lambda)}{(p-1)}\sum_{r=0}^{p-2}(-p)^{-\left\lfloor\frac{-r}{p-1}\right\rfloor-\left\lfloor\frac{-r}{p-1}\right\rfloor-\left\lfloor\frac{4r}{p-1}\right\rfloor+\left\lfloor\frac{2r}{p-1}\right\rfloor}
\frac{ \Gamma_p\left(\left\langle\frac{-r}{p-1}\right\rangle\right) \Gamma_p\left(\left\langle\frac{-r}{p-1}\right\rangle\right)\Gamma_p\left(\left\langle\frac{4r}{p-1}\right\rangle\right)\overline{\omega}^{r}(\frac{1}{2^6\lambda^2})}{\Gamma_p\left(\left\langle\frac{2r}{p-1}\right\rangle\right)}.
\end{align}
Now, using the product formula given in \eqref{prod-3} we have 
\begin{align}\label{eqn-1007}
   \Gamma_p\left(\left\langle\frac{4r}{p-1}\right\rangle\right)=\frac{\Gamma_p\left(\left\langle\frac{r}{p-1}\right\rangle\right)\Gamma_p\left(\left\langle\frac{r}{p-1}+\frac{1}{4}\right\rangle\right)\Gamma_p\left(\left\langle\frac{r}{p-1}+\frac{1}{2}\right\rangle\right)\Gamma_p\left(\left\langle\frac{r}{p-1}+\frac{3}{4}\right\rangle\right) {\overline\omega^r(4^4)}}{\Gamma_p\left(\left\langle\frac{1}{4}\right\rangle\right)\Gamma_p\left(\left\langle\frac{1}{2}\right\rangle\right)\Gamma_p\left(\left\langle\frac{3}{4}\right\rangle\right)} 
\end{align}
and
\begin{align}\label{eqn-1008}
   \Gamma_p\left(\left\langle\frac{2r}{p-1}\right\rangle\right)=\frac{\Gamma_p\left(\left\langle\frac{r}{p-1}\right\rangle\right)\Gamma_p\left(\left\langle\frac{r}{p-1}+\frac{1}{2}\right\rangle\right) {\overline\omega^r(4)}}{\Gamma_p\left(\left\langle\frac{1}{2}\right\rangle\right)}.
\end{align}
Furthermore, it is easy to see that for $0\leq r\leq p-2$ we have
\begin{equation}\label{eqn-1009}
    \left\lfloor\frac{4r}{p-1}\right\rfloor-\left\lfloor\frac{2r}{p-1}\right\rfloor=\left\lfloor\frac{1}{4}+\frac{r}{p-1}\right\rfloor+\left\lfloor\frac{3}{4}+\frac{r}{p-1}\right\rfloor.
\end{equation}
Substituting \eqref{eqn-1007} , \eqref{eqn-1008} and \eqref{eqn-1009} into \eqref{eqn-1010} and taking the transformation $r\rightarrow-r$ we deduce that
\begin{align}\label{eqn-1011}
A_2=-p\cdot  \phi(2\lambda)\cdot {_2G_2}\left[\begin{array}{cc}
 \frac{1}{4}    & \frac{3}{4} \\
  0    & 0  \end{array}|{\lambda^2}\right]_p.
\end{align}
Finally, substituting the values of $A_1$ and $A_2$ into \eqref{eqn-1004} we have the first equality. To obtain the second equality, we make the transformation $r\rightarrow -r-\frac{p-1}{2}$ in \eqref{eqn-1006} and proceeding similar steps as in the proof of the first equality we complete the proof.
\end{proof}

\end{document}
